\documentclass[10pt,a4paper]{amsart}
\usepackage[margin=19mm]{geometry}
\usepackage{amsmath,amssymb,mathtools}
\usepackage[scr=rsfs,cal=euler]{mathalfa}
\usepackage{xfrac}
\usepackage[unicode,hidelinks,bookmarks=false]{hyperref}
\newcommand{\ie}{i.e.\ }
\newcommand{\cf}{cf.\ }
\newcommand{\resp}{respectively\ }
\newcommand{\Subsection}{\subsection}
\providecommand{\nobreakdash}{\nobreak\hbox{-}}
\setlength{\emergencystretch}{2em}
\multlinegap0pt
\usepackage[shortlabels]{enumitem}
\usepackage{amssymb}
\usepackage{tikz}
\usepackage{mathtools}
\newtheorem{thm}{Theorem}
\newtheorem{prop}{Proposition}
\newcommand{\nat}{\mathbb{N}}
\title{\vspace*{-1.5cm} About Wess-Zumino-Witten equation \\ and Harder-Narasimhan potentials}
\author{Siarhei Finski}
\date{Source: arXiv:2407.06034v3 (CC BY 4.0)}
\begin{document}
\maketitle

\noindent\emph{Source and licence.} \vspace*{-1.5cm} About Wess-Zumino-Witten equation \\ and Harder-Narasimhan potentials. Version arXiv:2407.06034v3 (CC BY 4.0), distributed by arXiv under the Creative Commons Attribution 4.0 International licence, http://creativecommons.org/licenses/by/4.0/, as recorded in the arXiv licence metadata for this exact version and shown on the abstract page https://arxiv.org/abs/2407.06034v3. Full licence notice: \texttt{licenses/arxiv-2407.06034-CC-BY-4.0.txt}.\par\smallskip

\noindent\textit{Editorial scope.} Counted unit: the article's thm thm\_main2 (source0.tex lines 600--612). The article's complete original proof of it is delivered with it (source0.tex lines 1161--1230, one \texttt{proof} environment of 70 lines, which the article places away from the statement). No mathematics is written, repaired or re-derived: the statement and the proof are the article's own lines, in the article's own notation, and no retained line is substituted or re-flowed.  Retained uncounted as the source-local context the proof needs: lines 820--824; lines 1091--1093; lines 1112--1124; lines 1114--1123; lines 1130--1137.  Nothing was omitted from any retained span, so the package carries no omission record.  No retained line contains a citation, so the wrapper carries no bibliography and no bibliography entry.  No retained line contains a figure environment, an image-inclusion command or drawing code, so no image is delivered and the package declares no asset dependency.  Editorial formatting only, in addition: an AMS \texttt{amsart} preamble that restates the article's own declarations and macros used by the retained lines, namely the environments \texttt{thm}, \texttt{prop} on separate counters, together with the standard packages those lines need.  The retained environments print in this document as Theorem 1, Proposition 1 in order of appearance, whereas the article prints the same items with the numbering of its own full document.  The complete native source of the article as distributed in the arXiv e-print for this exact version is bundled unchanged as \texttt{sources/161978/source0.tex}.
	\begin{equation}\label{eq_geod_ray_bnd}
		\sup_{x \in X} \big| \phi(h^L, \mathcal{F})(x) \big|
		\leq
		\| \mathcal{F} \|.
	\end{equation}

	\begin{equation}\label{eq_resol_hn_filt}
		\mu_k^* E_k = \tilde{\mu}_k^* \mathcal{F}^{HN, k}_{\lambda_1} \supset \tilde{\mu}_k^* \mathcal{F}^{HN, k}_{\lambda_2} \supset \cdots \supset \tilde{\mu}_k^* \mathcal{F}^{HN, k}_{\lambda_{q_k}},
	\end{equation}

	\begin{thm}\label{thm_expl_const1}
		For any $\epsilon > 0$, there are $\delta > 0$, $k \in \nat$, such that for any $\delta$-approximate critical Hermitian structure $H_{\delta, k}$ on $E_k$, there is $s \in [0, + \infty[$, such that for the geodesic ray of Hermitian metrics $H_{\delta, k, s}$ on $\mu_k^* E_k$, departing from $\mu_k^* H_{\delta, k}$ and associated with the resolution of the Harder-Narasimhan filtration (\ref{eq_resol_hn_filt}), the $(1, 1)$-form $\omega_{\delta, k, s} := c_1(p_k^* L, FS(H_{\delta, k, s})^{\frac{1}{k}})$ verifies
		\begin{equation}\label{eq_expl_const1}
			\int_{X_k}
			\Big|
			\omega_{\delta, k, s}^{n + 1} \wedge \pi_k^* \mu_k^* \omega_B^{m - 1}
			-
			p_k^* {\rm{HN}}(\omega_{\delta, k, s}) \cdot \omega_{\delta, k, s}^n \wedge  \pi_k^* \mu_k^* \omega_B^m \cdot (n + 1)
			\Big|
			\leq
			\epsilon.
		\end{equation}
	\end{thm}

		\begin{equation}
			\int_{X_k}
			\Big|
			\omega_{\delta, k, s}^{n + 1} \wedge \pi_k^* \mu_k^* \omega_B^{m - 1}
			-
			p_k^* {\rm{HN}}(\omega_{\delta, k, s}) \cdot \omega_{\delta, k, s}^n \wedge  \pi_k^* \mu_k^* \omega_B^m \cdot (n + 1)
			\Big|
			\leq
			\epsilon.
		\end{equation}

	\begin{prop}\label{prop_bir_mod}
		For any $\epsilon > 0$, and a relatively Kähler $(1, 1)$-form $\omega_0$ on $X_0$ in the class $p_0^* c_1(L)$, there is a compact subset $K \subset B \setminus {\rm{Crit}}(\mu_0)$, such that for any compact subset $K' \subset B \setminus {\rm{Crit}}(\mu_0)$, verifying $K \subset K'$, there is a relatively Kähler $(1, 1)$-form $\omega$ on $X$ in the class $c_1(L)$, so that over $\pi_0^{-1}(\mu_0^{-1}(K'))$, the forms $\omega_0$ and $p_0^* \omega$ coincide, and 
		\begin{equation}\label{eq_prop_bir_mod}
			\int_{X_0 \setminus \pi_0^{-1}(\mu_0^{-1}(K'))} \big| \omega_0^{n + 1} \wedge \pi_0^* \mu_0^* \omega_B^{m - 1} \big| < \epsilon, 
			\qquad 
			\int_{X \setminus \pi^{-1}(K')} \big| \omega^{n + 1} \wedge \pi^* \omega_B^{m - 1} \big| < \epsilon.
		\end{equation}
	\end{prop}

	\begin{thm}\label{thm_main2}
		For any $\epsilon > 0$, there is a relatively Kähler $(1,1)$-form $\omega_{\epsilon} \in c_1(L)$ over $X$, verifying 
		\begin{equation}\label{eq_thm_main2}
			\int_X
			\Big|
			\omega_{\epsilon}^{n + 1} \wedge \pi^* \omega_B^{m - 1}
			-
			{\rm{HN}}(\omega_{\epsilon}) \cdot \omega_{\epsilon}^n \wedge \pi^* \omega_B^m \cdot (n + 1)
			\Big|
			\leq
			\epsilon.
		\end{equation}
	\end{thm}

	\begin{proof}[Proof of Theorem \ref{thm_main2}]
		We fix $\epsilon > 0$ and consider $\omega_{\delta, k, s}$ given by Theorem \ref{thm_expl_const1}.
		We fix a compact subset $K \subset B \setminus {\rm{Crit}}(\mu_k)$, given by Proposition \ref{prop_bir_mod}.
		We take a compact subset $K' \subset B \setminus {\rm{Crit}}(\mu_0)$, $K \subset K'$, which we specify later.
		Using Proposition \ref{prop_bir_mod}, we find a relatively Kähler $(1, 1)$-form $\omega_{\epsilon}$ on $X$ in the class $c_1(L)$, so that over $\pi_k^{-1} (\mu_k^{-1}(K'))$, the forms $\omega_{\delta, k, s}$ and $p_k^* \omega_{\epsilon}$ coincide, and
		\begin{equation}\label{eq_eps_away_from_soln}
			\int_{X_k \setminus \pi_k^{-1} (\mu_k^{-1}(K'))} \big| \omega_{\delta, k, s}^{n + 1} \wedge \pi_k^* \mu_k^* \omega_B^{m - 1} \big| < \epsilon, 
			\qquad 
			\int_{X \setminus \pi^{-1} (K')} \big| \omega_{\epsilon}^{n + 1} \wedge \pi^* \omega_B^{m - 1} \big| < \epsilon.
		\end{equation}
		Remark that due to coincidence of $\omega_{\delta, k, s}$ with $p_k^* \omega_{\epsilon}$ over $\pi_k^{-1} (\mu_k^{-1}(K'))$, we have $p_k^* {\rm{HN}}(\omega_{\delta, k, s}) = p_k^*( {\rm{HN}}(\omega_{\epsilon}))$ over $\pi_k^{-1} (\mu_k^{-1}(K'))$.
		From this and (\ref{eq_expl_const1}), the following bound is then immediate 
		\begin{equation}\label{eq_eps_away_from_soln2}
		\begin{aligned}
			\int_{X}
			\Big|
			\omega_{\epsilon}^{n + 1} \wedge & \pi^* \omega_B^{m - 1}
			-
			{\rm{HN}}(\omega_{\epsilon}) \cdot \omega_{\epsilon}^n \wedge  \pi^* \omega_B^m \cdot (n + 1)
			\Big|
			\\
			&
			\leq
			\int_{X_k \setminus \pi_k^{-1} (\mu_k^{-1}(K'))} \big| \omega_{\delta, k, s}^{n + 1} \wedge \pi_k^* \mu_k^* \omega_B^{m - 1} \big|
			+
			\int_{X \setminus \pi^{-1} (K')} \big| \omega_{\epsilon}^{n + 1} \wedge \pi^* \omega_B^{m - 1} \big|
			\\
			&
			{\phantom{\leq}}
			+
			(n + 1)
			\cdot
			\int_{X_k \setminus \pi_k^{-1} (\mu_k^{-1}(K'))}
			\Big|
			p_k^* {\rm{HN}}(\omega_{\delta, k, s}) \cdot \omega_{\delta, k, s}^n \wedge  \pi_k^* \mu_k^* \omega_B^m
			\Big|
			\\
			&
			{\phantom{\leq}}
			+
			(n + 1)
			\cdot
			\int_{X \setminus \pi^{-1} (K')}
			\Big|
			{\rm{HN}}(\omega_{\epsilon}) \cdot \omega_{\epsilon}^n \wedge  \pi^* \omega_B^m
			\Big| + \epsilon.
		\end{aligned}
		\end{equation}
		Remark, however, that by (\ref{eq_geod_ray_bnd}), there is $C > 0$, such that $|{\rm{HN}}(\omega_{\delta, k, s})| < C$, $|{\rm{HN}}(\omega_{\epsilon})| < C$.
		From this, (\ref{eq_eps_away_from_soln}) and (\ref{eq_eps_away_from_soln2}), we conclude that 
		\begin{multline}\label{eq_eps_away_from_soln3}
			\int_{X}
			\Big|
			\omega_{\epsilon}^{n + 1} \wedge \pi^* \omega_B^{m - 1}
			-
			{\rm{HN}}(\omega_{\epsilon}) \cdot \omega_{\epsilon}^n \wedge  \pi^* \omega_B^m \cdot (n + 1)
			\Big|
			\\
			\leq
			2 C \cdot (n + 1) \cdot
			\pi_* (c_1(L)^n)
			\cdot
			\int_{B \setminus K'} \omega_B^m 
			+
			3 \epsilon.
		\end{multline}
		The right-hand side of (\ref{eq_eps_away_from_soln3}) can be made smaller than $4 \epsilon$ by taking bigger $K'$.
		In particular, the forms $\omega_{\epsilon}$ satisfy Theorem \ref{thm_main2}, for $\epsilon := 4 \epsilon$.
		But as $\epsilon > 0$ was chosen in an arbitrary way, this finishes the proof.
	\end{proof}

\end{document}
